\section{Adaptación a dominios abiertos mediante rúbricas}

En los dominios abiertos, una misma consulta puede admitir múltiples respuestas válidas, cuyo valor depende del contexto, de la información incluida y de cómo se relacionan los conceptos. A diferencia de las tareas con una solución verificable, no suele existir una regla única que permita determinar si una respuesta es correcta en su conjunto. Aun así, algunos aspectos pueden verificarse, como la presencia de información relevante, el cumplimiento de ciertas restricciones o la inclusión de conceptos esperados, lo que ofrece un punto de partida para el entrenamiento. Para comprobar estos aspectos, es necesario reconocer que una misma idea puede expresarse de distintas maneras y que la presencia de un término no garantiza su uso adecuado en el contexto.

Para abordar esta dificultad, la metodología parte de una adaptación mediante ajuste fino supervisado (SFT) y utiliza respuestas de referencia de una partición reservada para extraer automáticamente las rúbricas, que reúnen los conceptos relevantes para cada consulta. La evaluación combina comprobaciones mediante código con un modelo de lenguaje que actúa como juez semántico, examinando los conceptos no detectados por esas comprobaciones para reconocer formulaciones equivalentes. A partir de esta evaluación se construye la recompensa que orienta el aprendizaje por refuerzo. Durante el entrenamiento, parte de los conceptos se presenta al modelo como ayuda y se retira progresivamente, mientras la rúbrica completa se conserva para evaluar sus respuestas. La figura~\ref{fig:intro-pipeline} resume este procedimiento.

\begin{figure}[H]
\centering
\begingroup
\singlespacing
\shorthandoff{<>}
\definecolor{pipelineink}{HTML}{253F53}
\definecolor{pipelineline}{HTML}{658AA8}
\definecolor{pipelinepaper}{HTML}{F7FAFC}
\definecolor{pipelineblue}{HTML}{C7DCEC}
\definecolor{pipelinepale}{HTML}{DBE7F0}
\definecolor{pipelineshadow}{HTML}{D6E0E8}
% Tarjeta vectorial con cabecera, cuerpo y sombra discreta.
\newcommand{\pipelinecard}[8]{%
 \fill[pipelineshadow,rounded corners=7pt]
 ({#2-#4/2-.055},{#3-#5/2-.075}) rectangle
 ({#2+#4/2-.055},{#3+#5/2-.075});
 \node[draw=pipelineline,fill=pipelinepaper,rounded corners=7pt,
   line width=.85pt,minimum width=#4cm,minimum height=#5cm,inner sep=0pt]
   (#1) at (#2,#3) {};
 \draw[draw=pipelineline,fill=#8,rounded corners=5pt,line width=.7pt]
 ({#2-#4/2+.13},{#3+#5/2-.55}) rectangle
 ({#2+#4/2-.13},{#3+#5/2-.13});
 \node[font=\sffamily\footnotesize\bfseries,text=pipelineink,align=center,inner sep=0pt]
 at (#2,{#3+#5/2-.34}) {#6};
 \node[font=\sffamily\small,text=pipelineink,align=center,inner sep=0pt]
 at (#2,{#3-.27}) {#7};
}
\resizebox{\textwidth}{!}{%
\begin{tikzpicture}[>=Stealth,
 flow/.style={->,draw=pipelineline,line width=1pt,rounded corners=4pt,line join=round},
 label/.style={font=\sffamily\footnotesize,fill=white,inner sep=3pt,text=pipelineink}]
\pipelinecard{data}{0}{0}{2.3}{1.95}{Consultas y}{respuestas\\del dominio}{pipelinepale}
\pipelinecard{sft}{2.9}{1.8}{2.1}{1.45}{Datos}{para SFT}{pipelineblue}
\pipelinecard{reserve}{2.9}{-1.8}{2.1}{1.95}{Datos}{reservados\\para GRPO}{pipelineblue}
\pipelinecard{model}{5.9}{1.8}{3.05}{1.45}{Entrenamiento}{con SFT}{pipelineblue}
\pipelinecard{extractor}{5.9}{-1.8}{2.4}{1.95}{LLM}{Extracción de\\criterios}{pipelinepale}
\pipelinecard{rubric}{8.9}{-1.8}{2.25}{1.45}{Rúbricas}{extraídas}{pipelinepale}
\fill[pipelineshadow,rounded corners=7pt] (10.62,-2.6) rectangle (13.07,-1.15);
\node[draw=pipelineline,fill=pipelinepaper,rounded corners=7pt,line width=.85pt,
 minimum width=2.45cm,minimum height=1.45cm,inner sep=0pt] (scaffold) at (11.9,-1.8) {};
\draw[draw=pipelineline,fill=pipelineblue,rounded corners=5pt,line width=.7pt]
 (10.805,-2.01) rectangle (12.995,-1.59);
\node[font=\sffamily\footnotesize\bfseries,text=pipelineink,inner sep=0pt]
 at (11.9,-1.8) {Scaffolding};
\pipelinecard{grpo}{15.1}{0}{3.05}{1.95}{Entrenamiento}{con GRPO\\con rúbricas}{pipelineblue}
\pipelinecard{final}{18.3}{0}{2.1}{1.45}{Modelo}{Final}{pipelineblue}
\draw[flow] (data.east) -- (1.5,0) |- (sft.west);
\draw[flow] (1.5,0) |- (reserve.west);
\draw[flow] (sft.east) -- (model.west);
\draw[flow] (reserve.east) -- (extractor.west);
\draw[flow] (extractor.east) -- (rubric.west);
\draw[flow] (reserve.north) -- (2.9,-.25) --
 node[label,above]{Consultas} (11.9,-.25) -- (scaffold.north);
\draw[flow] (rubric.east) -- (scaffold.west);
\draw[flow] (model.east) -| (grpo.north);
\draw[flow] (scaffold.east) -- (13.35,-1.8) |- (grpo.west);
\draw[flow] (rubric.south) -- (8.9,-3.2) --
 node[label,above]{Rúbrica completa} (15.1,-3.2) -- (grpo.south);
\draw[flow] (grpo.east) -- (final.west);
\end{tikzpicture}%
}
\endgroup

\caption{Metodología de adaptación mediante SFT y GRPO, con rúbricas extraídas de respuestas de referencia, evaluación híbrida y retirada progresiva de ayudas.}
\label{fig:intro-pipeline}
\end{figure}

Los trabajos previos abordan distintos componentes de este procedimiento. Rubicon estudia las rúbricas como señal de recompensa~\cite{huang2025rubicon}, mientras que ARES automatiza la construcción de preguntas, respuestas y criterios a partir de documentos~\cite{li2026ares} y RuscaRL incorpora ayudas basadas en rúbricas que se retiran progresivamente~\cite{zhou2026ruscarl}. Frente a estos antecedentes, el presente trabajo aprovecha pares de consulta y respuesta existentes para realizar la adaptación supervisada y extraer criterios de una partición reservada, independiente del benchmark externo. Su aporte se concentra en integrar la evaluación híbrida y la retirada progresiva de ayudas en ese procedimiento de adaptación de dominio.

El dominio clínico ofrece un caso de estudio relevante para esta metodología, porque sus consultas requieren respuestas abiertas que combinan conocimientos especializados con información propia de cada situación. Evaluarlas exige considerar tanto los conceptos incluidos como las posibles omisiones y la adecuación de las recomendaciones al contexto~\cite{arora2025healthbench}. Por ello, las consultas clínicas permiten estudiar si la recompensa favorece respuestas que recuperan información relevante y la utilizan adecuadamente en el contexto del paciente.

Para estudiar los efectos del entrenamiento, la evaluación compara el modelo adaptado mediante SFT con su posterior optimización mediante aprendizaje por refuerzo guiado por rúbricas. Para ello, se utiliza HealthBench como evaluación externa, con rúbricas independientes de las empleadas durante el entrenamiento~\cite{arora2025healthbench}. El análisis examina los cambios en distintas dimensiones de las respuestas y se complementa con ejemplos que permiten identificar mejoras, regresiones y limitaciones de la propuesta.


La investigación se desarrolla en dos etapas. En la primera, los experimentos con Qwen2.5-7B permitieron evaluar la integración inicial de SFT, rúbricas y GRPO e identificar mejoras en algunas dimensiones junto con regresiones y limitaciones del evaluador. A partir de esos aprendizajes, la segunda etapa utiliza preguntas y respuestas procedentes de Rubrics as Rewards (RaR) para construir nuevas rúbricas y estudiar distintas formas de puntuarlas~\cite{gunjal2025rar}. Esta etapa trabaja con Qwen2.5-3B y otro diseño experimental; su relación con la evaluación inicial es metodológica, como se precisa en la sección~\ref{sec:resultados-modelos}.
